% Commutative diagrams (tikz-cd) between and inside shaded and leftbar boxes
% (framed.sty), and a framed box whose frame is an \fcolorbox: every framed
% environment measures its frame with \fb@sizeofframe. Written for the corpus
% (lane P6-INFDESC-PAGE, BOX-MEMO's sweep coverage).
\documentclass[11pt]{article}
\usepackage[a5paper,margin=16mm]{geometry}
\usepackage{amsmath,amssymb}
\usepackage{xcolor}
\usepackage{tikz-cd}
\usepackage{framed}
\definecolor{shadecolor}{RGB}{240,240,232}
\newenvironment{note}{%
  \def\FrameCommand{\fboxsep=4pt\fcolorbox{gray}{white}}%
  \MakeFramed{\advance\hsize-\width \FrameRestore}}%
  {\endMakeFramed}

\begin{document}

\section*{Quotients}

A surjection $p : X \to Z$ induces an equivalence relation $\sim$ on $X$, and
$p$ factors through the quotient map:
\begin{center}
\begin{tikzcd}
X \arrow[r, "q"] \arrow[dr, "p"'] & X/{\sim} \arrow[d, "\bar p", "\cong"'] \\
& Z
\end{tikzcd}
\end{center}

\begin{shaded}
\textbf{Universal property.} For every function $f : X \to Y$ that is constant
on equivalence classes there is exactly one $\bar f$ with $f = \bar f \circ q$:
\[
\begin{tikzcd}
X \arrow[r, "q"] \arrow[rd, "f"'] & X/{\sim} \arrow[d, dashed, "\exists!\,\bar f"] \\
& Y
\end{tikzcd}
\]
\end{shaded}

\begin{leftbar}
The three descriptions are equivalent: equivalence relations on $X$, partitions
of $X$ and surjections with domain $X$ (up to bijection of the codomain).
\begin{center}
\begin{tikzcd}[row sep=large, column sep=small]
& \text{relations} \arrow[dl, leftrightarrow] \arrow[dr, leftrightarrow] & \\
\text{partitions} \arrow[rr, leftrightarrow, dashed] & & \text{surjections}
\end{tikzcd}
\end{center}
\end{leftbar}

\begin{note}
\textbf{Exercise.} Show that the composite of two quotient maps is a quotient
map, and draw the square
\[
\begin{tikzcd}
A \arrow[r] \arrow[d] & B \arrow[d] \\
C \arrow[r] & D
\end{tikzcd}
\]
for the quotients of $\mathbb{Z}$ by $4\mathbb{Z}$ and $2\mathbb{Z}$.
\end{note}

\begin{shaded}
\textbf{Products.} The product $X \times Y$ comes with projections, and maps
into it are pairs of maps:
\[
\begin{tikzcd}
& W \arrow[dl, "f"'] \arrow[dr, "g"] \arrow[d, dashed, "{(f,g)}"] & \\
X & X \times Y \arrow[l, "\pi_1"] \arrow[r, "\pi_2"'] & Y
\end{tikzcd}
\]
Text after the diagram fills the box so that it continues to the next page if
the page is short: the frame is then measured and drawn for each part. A box
split across pages is a framed environment's hardest case, and the one
\texttt{framed} measures most often.
\end{shaded}

\begin{note}
\textbf{Coproducts.} Dually, maps out of a disjoint union are pairs of maps.
\[
\begin{tikzcd}
X \arrow[r, "\iota_1"] \arrow[dr, "f"'] & X \sqcup Y \arrow[d, dashed] & Y \arrow[l, "\iota_2"'] \arrow[dl, "g"] \\
& W &
\end{tikzcd}
\]
\end{note}

\end{document}
